\documentclass[a4paper]{article}
\usepackage{titlesec}
\usepackage{geometry}
\geometry{hmargin={2cm,2cm}, vmargin={3.5cm,3.5cm}}
\usepackage{graphicx}
\usepackage[table]{xcolor}
\usepackage{hyperref}
\hypersetup{pdfborder = {0 0 0},
colorlinks,
citecolor=black,
filecolor=black,
linkcolor=black,
urlcolor=blue}
\definecolor{logo}{RGB}{20,99,163}
\usepackage[T1]{fontenc}
\usepackage[utf8]{inputenc}
\usepackage[english,italian]{babel}
\usepackage{fancyhdr}
\usepackage{lastpage}
\usepackage{comment}

% creazione comando subsubsubsection
\titleclass{\subsubsubsection}{straight}[\subsection]

\newcounter{subsubsubsection}[subsubsection]
\renewcommand\thesubsubsubsection{\thesubsubsection.\arabic{subsubsubsection}}
\titleformat{\subsubsubsection}
    {\normalfont\normalsize\bfseries}{\thesubsubsubsection}{1em}{}
    \titlespacing*{\subsubsubsection}
    {0pt}{3.25ex plus 1ex minus .2ex}{1.5ex plus .2ex}

\makeatletter
\def\toclevel@subsubsubsection{4}
\def\l@subsubsubsection{\@dottedtocline{4}{7em}{4em}}
\makeatother

\setcounter{secnumdepth}{4}
\setcounter{tocdepth}{4}
%fine creazione comando


%Comandi di impaginazione uguale per tutti i documenti
\pagestyle{fancy}
\lhead{\includegraphics[scale=0.2]{../../../../Images/logo.png}}
%Titolo del documento
\rhead{\large Verbale Esterno del 2021-08-03}
%\rfoot{\thepage}
\cfoot{Pagina \thepage{} di \pageref{LastPage}}
\setlength{\headheight}{45pt}
\renewcommand{\footrulewidth}{0.4pt}
\begin{document}

\begin{titlepage}
    \begin{center}
        \includegraphics{../../../../Images/logo.png}\\
        \vspace{20px}
        \textcolor{logo}{\hrulefill}\\
        \vspace{20px}
        \textbf{\huge\textcolor{logo}{Verbale Esterno del 2021-08-03}}\\
        \vspace{10px}
        \textcolor{logo}{\hrulefill}\\
        \vspace{40px}
        \textbf{\Large Informazioni sul documento}\\
        \vspace{20px}
        \begin{tabular}{p{100px} | p{100px}}
            \textbf{Versione}     & 1.0.0                                                                      \\
            \textbf{Responsabile} & Marco Barbaresco                                                           \\
            \textbf{Redattore}    & Nicolò Giaccone                                                            \\
            \textbf{Verificatore} & Manuel Varo                                                                \\
            \textbf{Destinatari}  & Prof. Tullio Vardanega \newline Prof. Riccardo Cardin, \newline TechSWEave \\
            \textbf{Stato}        & Approvato                                                                  \\
            \textbf{Uso}          & Esterno                                                                    \\
        \end{tabular}\\
        \vspace{60px}
        \href{mailto:techsweave@gmail.com}{techsweave@gmail.com}\\

    \end{center}
\end{titlepage}

%INDICE
\newpage
\tableofcontents
\newpage

\section{Informazioni riunione}
\begin{itemize}
    \item Data e ora: 2021-08-03 ore 16:00-16:30.
    \item Luogo: Google Meet fornito dai proponenti.
    \item Partecipanti interni: Samuele De Simone, Manuel Varo, Tito Scutari, Nicolò Giaccone, Amedeo Meggiolaro, Marco Barbaresco.
    \item Partecipanti esterni: Milo Ertola, Alessandro Maccagnan.
    \item Segretario: Amedeo Meggiolaro.
\end{itemize}
\section{Argomenti affrontati}
Durante l'incontro è stato mostrato il prodotto ai proponenti, chiedendo loro pareri su criticità del prodotto da correggere e sul prodotto in genrale.
\section{Verbale della riunione}
\begin{itemize}
    \item Aggiornamento del proponente sullo stato del progetto;
    \item Revisione del prodotto, da cui sono emerse alcune criticità:
          \begin{itemize}
              \item un venditore deve poter accedere ai dettagli dell'utente dalla lista degli ordini;
              \item maggiore distinzione tra la PDP per il cliente e quella per il venditore, che al momento risultano essere troppo simili;
              \item rimuovere il multi loading dalla lista degli ordini.
          \end{itemize}
    \item Fissata una riunione da effettuare in seguito alla \textit{Revisione di Qualifica}, per discutere dell'andamento della revisione e per organizzare le attività da svolgere in vista della \textit{Revisione di Accettazione}.
\end{itemize}
\newpage
\section{Resoconto}
\begin{center}
    \textbf{\Large Registro delle decisioni}\\
    \vspace{10px}
    \begin{table}[h!]
        \centering
        \rowcolors{2}{logo!10}{logo!40}
        \renewcommand{\arraystretch}{1.8}
        \begin{tabular}{p{200px} p{200px}}
            \rowcolor{logo!70} \textbf{Codice} & \textbf{Decisione}                                                                \\
            VE\_2021\_08\_03.1                 & Aggiunto accesso al dettaglio dell'utente dalla lista degli ordini del venditore. \\
            VE\_2021\_08\_03.2                 & Diversificata la PDP del cliente rispetto a quella del venditore.                 \\
            VE\_2021\_08\_03.3                 & Rimosso il multi loading dalla lista degli ordini.                                \\
        \end{tabular}
    \end{table}
\end{center}

\end{document}